\documentclass[../../../include/open-logic-section]{subfiles}

\begin{document}

\olfileid{sth}{z}{pairs}
\olsection{جوڑیاں}

اگلا مسلّمہ جیہڑا ویکھنا اے، ایہ اے:

\begin{axiom}[جوڑیاں]
کسے وی دو سیٹاں $a, b$ لئی سیٹ $\{a, b\}$ موجود اے۔
\[
	\forall a \forall b \exists P \forall x (x \in P \liff (x = a \lor x = b))
\]
\end{axiom}

مرحلہ وار تصور نال ایس مسلّمے دا جواز ایویں دیندے آں۔ منّ لو $a$
مرحلے $S$ اُتے تے $b$ مرحلے $T$ اُتے موجود اے۔ منّ لو $M$، مرحلیاں
$S$ تے $T$ وچوں جیہڑا پچھوں آؤندا اے، اوہ اے۔ چونکہ $a$ تے $b$ دونیں
مرحلے $M$ اُتے موجود نیں، ایس لئی سیٹ $\{a,b\}$ اک ممکن اکٹھ اے جیہڑا
$M$ توں پچھوں آؤن والے کسے وی مرحلے اُتے دستیاب اے (جیہڑا وی وڈا ہووے)۔

پر ذرا ٹھہرو!{} اسیں کیوں منّیئے کہ $M$ توں پچھوں کوئی مرحلے
\emph{موجود نیں}؟ جے کوئی نہیں تاں ساڈا جواز نہیں چلے گا۔ سو جوڑیاں دا
جواز دین لئی سانوں \olref[sth][z][story]{sec} دے بیان وچ اک ہور
اصول شامل کرنا پئے گا:
\begin{enumerate}
	\item[] \stagessucc۔ کوئی آخری مرحلہ نہیں۔
\end{enumerate}
کی ایس اصول دا جواز اے؟ شون فیلڈ دے بیان وچ ایہ گل \emph{صاف لفظاں وچ}
نہیں کہی گئی سی کہ کوئی آخری مرحلہ نہیں۔ فیر وی، جے ایہ (ٹھیک ٹھیک
آکھیے تاں) ساڈے بیان وچ اک وادھا اے، ایہ ایس بنیادی خیال نال میل کھاندا
اے کہ سیٹ مرحلیاں وچ بندے نیں۔ اگے اسیں اینوں منّ لواں گے۔ سو جوڑیاں دا
مسلّمہ وی منّ لواں گے۔

ایس نویں مسلّمے نال اسیں کئی ہور سیٹاں دی موجودگی ثابت کر سکدے آں۔ مثلاً:

\begin{prop}\ollabel{prop:pairsconsequences}
کسے وی دو سیٹاں $a$ تے $b$ لئی ہِیٹھلے سیٹ موجود نیں:
	\begin{enumerate}
		\item\ollabel{singleton} $\{a\}$
		\item\ollabel{binunion} $a \cup b$
		\item\ollabel{tuples} $\tuple{a, b}$
	\end{enumerate}
\end{prop}

\begin{proof}
\olref{singleton}۔ جوڑیاں دے اصول نال $\{a, a\}$ موجود اے، تے عنصراں
نال تعیّن دے اصول نال ایہ $\{a\}$ اے۔

\olref{binunion}۔ جوڑیاں دے اصول نال $\{a, b\}$ موجود اے۔
ہُن $a \cup b = \bigcup \{a, b\}$ اتحاد دے اصول نال موجود اے۔

\olref{tuples}۔ \olref{singleton} نال $\{a\}$ موجود اے۔ جوڑیاں دے
اصول نال $\{a, b\}$ موجود اے۔ ہُن $\{\{a\}, \{a, b\}\} = \tuple{a, b}$
وی جوڑیاں دے اصول نال موجود اے۔
\end{proof}

\begin{prob}
وکھاؤ کہ کسے وی سیٹاں $a, b, c$ لئی سیٹ $\{a, b, c\}$ موجود اے۔
\end{prob}

\begin{prob}
وکھاؤ کہ کسے وی سیٹاں $a_1, \ldots, a_n$ لئی سیٹ $\{a_1, \ldots, a_n\}$
موجود اے۔
\end{prob}

%\begin{proof} By two applications of Pairs, $\{\{a_1, a_2\}, \{a_1,
%   a_3\}\}$ exists. By Union and Extensionality, $\bigcup \{\{a_1,
%   a_2\}, \{a_1, a_3\}\} = \{a_1, a_2, a_3\}$ exists. Repeat this
%   trick as often as necessary. \end{proof}

\end{document}
